% Page 044

\seite{44}

\abschnitt{Spaziergang}
\pstart
Am 12.\ September wurde mit allen vier\ob
Pfarreischulen ein Spaziergang nach\ob
Echsburg\unsicher{} unternommen.

\pend
\abschnitt{Ernte}
\pstart

Die Ernte war in diesem Jahre befriedigend.\ob
Das Wetter war gut, doch blieb das Obst\ob
gering. Die Kartoffeln erbrachten einen\ob
mittleren Ertrag, doch trat stellenweise\ob
die Fäule auf. Das Getreide wurde gut\ob
eingebracht.

\pend
\abschnitt{Weihnachtsfeier}
\pstart

Zum Geburtstag Seiner Majestät des Kaisers\ob
fand eine gemeinsame Feier statt. Die\ob
Kinder hielten Ansprachen, und Herr\ob
Pfarrer von Seffern hielt die Festrede.

Wegen Krankheit wurde die Schule vom\ob
15.\ Januar bis zum 3.\ März geschlossen.

\pend
\abschnitt{den 18. März 1908}
\pstart

Schmitt, Pfarrer. {[Unterschrift]}

\pend
\abschnitt{Osterprüfung}
\pstart

Am 18.\ März fand durch den Herrn\ob
Ortsschulinspektor Pfarrer Schmitt die\ob
übliche Osterprüfung statt.

\pend
\jahresueberschrift{Das Schuljahr 1908-09}
\pstart

Es wurden entlassen 5 Knaben und 4\ob
Mädchen. Die Knaben widmen sich der\ob
Landwirthschaft, die Mädchen der Haus\ohb
wirthschaft. Aufgenommen wurden 8 Knaben\ob
und 1 Mädchen. Alle Kinder sind katholisch\ob
und willfährig; ein Knabe ist nicht ganz\ob
willfährig.

\pend
\abschnitt{Spaziergang}
\pstart

Am 5.\ September wurde mit allen vier\ob
Pfarreischulen ein Spaziergang nach\ob
St.\ Thomas unternommen.
\pend
